\NormalFont\ShortTitle{2 Corinteni}

{\MT A doua scrisoare a lui Pavel către Corinteni


\par }\ChapOne{1}{\PP \VerseOne{1}Pavel, apostol al lui Isus Hristos, prin voia lui Dumnezeu, și Timotei, fratele nostru, către adunarea lui Dumnezeu care este în Corint, împreună cu toți sfinții care sunt în toată Ahaia:

\VS{2}Har și pace vouă, de la Dumnezeu, Tatăl nostru, și de la Domnul Isus Hristos!


\par }{\PP \VS{3}Binecuvântat să fie Dumnezeul și Tatăl Domnului nostru Isus Hristos, Tatăl îndurărilor și Dumnezeul oricărei mângâieri,

\VS{4}care ne mângâie în toate necazurile noastre, ca să putem mângâia și noi pe cei ce se află în vreun necaz, prin mângâierea cu care noi înșine suntem mângâiați de Dumnezeu.

\VS{5}Căci, după cum suferințele lui Hristos abundă pentru noi, tot așa și mângâierea noastră abundă prin Hristos.

\VS{6}Dar dacă suntem în suferință, este pentru mângâierea și mântuirea voastră. Dacă noi suntem mângâiați, este pentru mângâierea voastră, care produce în voi răbdarea răbdării acelorași suferințe pe care le suferim și noi.

\VS{7}Nădejdea noastră pentru voi este statornică, știind că, întrucât sunteți părtași la suferințe, sunteți și la mângâiere.


\par }{\PP \VS{8}Căci nu vrem, fraților, să nu fiți informați despre nenorocirea noastră, care ni s-a întâmplat în Asia, că am fost împovărați peste măsură de mult, mai presus de puterile noastre, încât am ajuns la disperare de viață.

\VS{9}Da, noi înșine am avut în noi înșine condamnarea la moarte, ca să nu ne încredem în noi înșine, ci în Dumnezeu, care învie morții,

\VS{10}care ne-a izbăvit dintr-o moarte atât de mare și care ne izbăvește și ne izbăvește, în care ne-am pus nădejdea că ne va izbăvi și pe noi în continuare,

\VS{11}ajutând și voi împreună în favoarea noastră prin rugămințile voastre, pentru ca, pentru darul care ne-a fost dat prin intermediul multora, să se aducă mulțumiri de către mulți în numele vostru.


\par }{\PP \VS{12}Căci lauda noastră este aceasta: mărturia conștiinței noastre că în sfințenie și în sinceritate față de Dumnezeu, nu în înțelepciune trupească, ci în harul lui Dumnezeu, ne-am purtat în lume și mai mult decât atât față de voi.

\VS{13}Căci nu vă scriem altceva decât ceea ce citiți sau chiar recunoașteți, și sper că veți recunoaște până la capăt —

\VS{14}așa cum și voi ne-ați recunoscut în parte — că noi suntem lauda voastră, așa cum și voi sunteți lauda noastră, în ziua Domnului nostru Isus.


\par }{\PP \VS{15}În această încredere, am hotărât să vin mai întâi la voi, ca să vă dau un al doilea folos,

\VS{16}să trec prin voi în Macedonia, iar din Macedonia să vin la voi și să fiu trimis de voi în Iudeea.

\VS{17}Așadar, când am plănuit aceasta, am dat oare dovadă de nestatornicie? Sau lucrurile pe care le plănuiesc, le plănuiesc eu după trup, ca la mine să fie “Da, da” și “Nu, nu”?”

\VS{18}Dar, cum Dumnezeu este credincios, cuvântul nostru față de voi nu a fost “Da și nu”.

\VS{19}Căci Fiul lui Dumnezeu, Isus Hristos, care a fost propovăduit printre voi de noi — de mine, de Silvanus și de Timotei — nu a fost “Da și nu”, ci în El este “Da”.

\VS{20}Căci oricât de multe ar fi promisiunile lui Dumnezeu, în el este “Da”. De aceea și prin el este “Amin”, pentru slava lui Dumnezeu prin noi.


\par }{\PP \VS{21}Și Cel ce ne-a întemeiat împreună cu voi în Hristos și ne-a uns, este Dumnezeu,

\VS{22}care ne-a pecetluit și ne-a dat în inimile noastre plata Duhului.


\par }{\PP \VS{23}Dar chem pe Dumnezeu ca martor al sufletului meu, că, pentru a vă cruța, n-am venit la Corint.

\VS{24}Noi nu stăpânim credința voastră, ci suntem împreună-lucrători cu voi pentru bucuria voastră. Căci voi rămâneți tari în credință.



\par }\Chap{2}{\PP \VerseOne{1}Dar am hotărât să nu mai vin la voi cu durere.

\VS{2}Căci, dacă vă întristez, cine mă va bucura pe mine, dacă nu cel care va fi întristat de mine?

\VS{3}Și v-am scris chiar acest lucru, pentru ca, atunci când voi veni, să nu am parte de întristare din partea celor de care ar trebui să mă bucur; având încredere în voi toți că bucuria mea va fi împărtășită de voi toți.

\VS{4}Pentru că din multă suferință și din angoasă sufletească v-am scris cu multe lacrimi, nu ca să vă întristez, ci ca să cunoașteți dragostea pe care o am din belșug pentru voi.


\par }{\PP \VS{5}Dar, dacă cineva a provocat întristare, nu mie mi-a provocat întristare, ci în parte, ca să nu vă presez prea mult, vouă tuturor.

\VS{6}Această pedeapsă care a fost aplicată de cei mulți este suficientă pentru unul ca acesta;

\VS{7}așa că, dimpotrivă, mai degrabă trebuie să-l iertați și să-l mângâiați, ca nu cumva unul ca acesta să fie înghițit de durerea lui excesivă.

\VS{8}De aceea vă rog să vă confirmați dragostea față de el.

\VS{9}Căci în acest scop am scris și eu, ca să cunosc dovada despre voi, dacă sunteți ascultători în toate lucrurile.

\VS{10}Acum și eu iert oricui iertați ceva. Căci, dacă într-adevăr am iertat ceva, am iertat aceluia pentru voi, în prezența lui Hristos,

\VS{11}pentru ca Satana să nu obțină niciun avantaj asupra noastră, căci nu suntem neștiutori de uneltirile lui.


\par }{\PP \VS{12}Când am venit la Troa pentru vestea lui Hristos și mi s-a deschis o ușă în Domnul,

\VS{13}nu am avut nici o ușurare pentru duhul meu, pentru că nu l-am găsit pe Tit, fratele meu; dar, luându-mi rămas bun de la ei, am plecat în Macedonia.


\par }{\PP \VS{14}Mulțumiri lui Dumnezeu, care ne conduce totdeauna în Hristos, și care, prin noi, descoperă în orice loc mireasma dulce a cunoștinței Lui.

\VS{15}Căci noi suntem o aromă dulce a lui Hristos pentru Dumnezeu, atât în cei ce se mântuiesc, cât și în cei ce pier:

\VS{16}pentru unii o duhoare din moarte în moarte, pentru ceilalți o aromă dulce din viață în viață. Cine este suficient pentru aceste lucruri?.

\VS{17}Căci noi nu suntem ca mulți, care colportăm cuvântul lui Dumnezeu. Ci ca din sinceritate, ci ca de la Dumnezeu, în fața lui Dumnezeu, noi vorbim în Hristos.



\par }\Chap{3}{\PP \VerseOne{1}Începem oare să ne lăudăm din nou? Sau avem nevoie, ca unii, de scrisori de laudă către voi sau de la voi?

\VS{2}Voi sunteți scrisoarea noastră, scrisă în inimile noastre, cunoscută și citită de toți oamenii,

\VS{3}fiind descoperit că sunteți o scrisoare a lui Hristos, slujită de noi, scrisă nu cu cerneală, ci cu Duhul Dumnezeului celui viu; nu pe table de piatră, ci pe table care sunt inimi de carne.


\par }{\PP \VS{4}O astfel de încredere avem în Hristos față de Dumnezeu, prin Hristos,

\VS{5}nu că ne-ar ajunge ceva de la noi înșine, ca să dăm socoteală de la noi înșine, ci că ne este de la Dumnezeu,

\VS{6}care ne-a făcut și pe noi destoinici, ca slujitori ai unui legământ nou, nu al literei, ci al Duhului. Căci litera ucide, dar Duhul dă viață.


\par }{\PP \VS{7}Dar dacă slujba morții, scrisă și gravată pe pietre, a venit cu slavă, încât copiii lui Israel nu puteau să privească cu stăruință la fața lui Moise, din pricina slavei feței lui, care era trecătoare,

\VS{8}nu va fi oare slujba Duhului cu mult mai multă slavă?

\VS{9}Căci dacă slujba condamnării are glorie, slujba dreptății depășește cu mult mai mult în glorie.

\VS{10}Căci, cu siguranță, ceea ce a fost făcut glorios nu a fost făcut glorios în această privință, din cauza gloriei care întrece.

\VS{11}Căci dacă ceea ce trece a fost cu glorie, cu atât mai mult ceea ce rămâne este în glorie.


\par }{\PP \VS{12}De aceea, având o astfel de nădejde, ne folosim de o mare îndrăzneală în vorbire,

\VS{13}și nu ca Moise, care și-a pus un văl pe față, ca să nu privească cu stăruință copiii lui Israel la sfârșitul a ceea ce era pe cale să treacă.

\VS{14}Dar mintea lor a fost împietrită, căci până în ziua de azi, la citirea vechiului legământ, același văl rămâne, pentru că în Cristos el trece.

\VS{15}Dar până în ziua de azi, când se citește Moise, un văl stă pe inima lor.

\VS{16}Dar de câte ori cineva se întoarce la Domnul, vălul este îndepărtat.

\VS{17}Or, Domnul este Duhul; și unde este Duhul Domnului, acolo este libertate.

\VS{18}Dar noi toți, văzând cu fața descoperită slava Domnului ca într-o oglindă, ne transformăm în același chip, din slavă în slavă, ca de la Domnul, Duhul.



\par }\Chap{4}{\PP \VerseOne{1}De aceea, având această slujbă, nu slăbim, pentru că am primit îndurare.

\VS{2}ci am renunțat la lucrurile ascunse ale rușinii, nu umblând cu viclenie și nici nu mânuind cu viclenie cuvântul lui Dumnezeu, ci prin manifestarea adevărului, recomandându-ne conștiinței oricărui om înaintea lui Dumnezeu.

\VS{3}Chiar dacă Buna Vestire a noastră este voalată, ea este voalată în cei muribunzi,

\VS{4}în care dumnezeul acestei lumi a orbit mințile celor necredincioși, pentru ca lumina Bunei Vestiri a slavei lui Cristos, care este chipul lui Dumnezeu, să nu răsară asupra lor.

\VS{5}Căci noi nu ne predicăm pe noi înșine, ci pe Cristos Isus ca Domn, și pe noi înșine ca slujitori ai voștri pentru Isus,

\VS{6}întrucât Dumnezeu, care a spus: “Lumina va străluci din întuneric”, este cel care a strălucit în inimile noastre pentru a da lumina cunoașterii gloriei lui Dumnezeu pe fața lui Isus Cristos.


\par }{\PP \VS{7}Dar noi avem această comoară în vase de lut, pentru ca puterea să fie de la Dumnezeu și nu de la noi înșine.

\VS{8}Suntem apăsați din toate părțile, dar nu zdrobiți; suntem în încurcătură, dar nu la disperare;

\VS{9}urmăriți, dar nu părăsiți; loviți, dar nu distruși;

\VS{10}purtând mereu în trup punerea la moarte a Domnului Isus, pentru ca viața lui Isus să se arate și în trupul nostru.

\VS{11}Căci noi, care trăim, suntem întotdeauna dați la moarte din pricina lui Isus, pentru ca și viața lui Isus să se arate în trupul nostru muritor.

\VS{12}Astfel, deci, moartea lucrează în noi, dar viața în voi.


\par }{\PP \VS{13}Dar având același duh de credință, după cum este scris: “Am crezut și de aceea am vorbit”. Și noi credem, și de aceea și vorbim,

\VS{14}știind că Cel care l-a înviat pe Domnul Isus ne va învia și pe noi cu Isus și ne va prezenta împreună cu voi.

\VS{15}Căci toate sunt pentru voi, pentru ca harul, înmulțindu-se prin cei mulți, să facă ca mulțumirile să abunde, spre slava lui Dumnezeu.


\par }{\PP \VS{16}De aceea, nu leșinăm, ci, deși exteriorul nostru se strică, lăuntrul nostru se înnoiește în fiecare zi.

\VS{17}Căci necazul nostru ușor, care este de moment, ne lucrează din ce în ce mai mult o greutate veșnică de glorie,

\VS{18}în timp ce noi nu ne uităm la lucrurile care se văd, ci la cele care nu se văd. Căci lucrurile care se văd sunt vremelnice, dar lucrurile care nu se văd sunt veșnice.



\par }\Chap{5}{\PP \VerseOne{1}Căci știm că, dacă se va desființa cortul nostru pământesc, avem o clădire de la Dumnezeu, o casă nefăcută de mână, veșnică, în ceruri.

\VS{2}Căci, cu siguranță, în aceasta gemem, dorind să fim îmbrăcați cu locuința noastră care este din ceruri,

\VS{3}dacă, într-adevăr, fiind îmbrăcați, nu vom fi găsiți goi.

\VS{4}Căci, într-adevăr, noi, care suntem în acest cort, gemem, fiind împovărați, nu că am dori să fim dezbrăcați, ci că dorim să fim îmbrăcați, pentru ca ceea ce este muritor să fie înghițit de viață.

\VS{5}Or, cel care ne-a făcut tocmai pentru acest lucru este Dumnezeu, care ne-a dat și avansul Duhului.


\par }{\PP \VS{6}De aceea suntem totdeauna încrezători și știm că, dacă suntem acasă în trup, suntem departe de Domnul,

\VS{7}căci umblăm prin credință, nu prin vedere.

\VS{8}Suntem curajoși, zic, și suntem dispuși mai degrabă să fim absenți din trup și să fim acasă la Domnul.

\VS{9}De aceea și noi ne propunem, fie că suntem acasă, fie că suntem absenți, să îi fim binevoitori.

\VS{10}Căci toți trebuie să fim descoperiți înaintea scaunului de judecată al lui Hristos, pentru ca fiecare să primească cele din trup, după cum a făcut, fie că este bun, fie că este rău.


\par }{\PP \VS{11}Cunoscând deci frica de Domnul, noi convingem pe oameni, dar suntem descoperiți lui Dumnezeu, și sper că suntem descoperiți și în conștiințele voastre.

\VS{12}Căci nu ne recomandăm din nou vouă, ci vorbim ca să vă dăm prilej de laudă în numele nostru, ca să aveți ce răspunde celor care se laudă în aparență și nu în inimă.

\VS{13}Căci, dacă suntem în afară de noi înșine, este pentru Dumnezeu. Sau dacă suntem cu mintea trează, este pentru voi.

\VS{14}Căci dragostea lui Hristos ne constrânge, pentru că noi judecăm astfel: că unul a murit pentru toți, deci toți au murit.

\VS{15}El a murit pentru toți, pentru ca cei care trăiesc să nu mai trăiască pentru ei înșiși, ci pentru cel care a murit și a înviat pentru ei.


\par }{\PP \VS{16}De aceea nu mai cunoaștem pe nimeni după trup de acum înainte. Chiar dacă am cunoscut pe Hristos după trup, acum nu-L mai cunoaștem astfel.

\VS{17}Prin urmare, dacă cineva este în Hristos, este o făptură nouă. Lucrurile vechi au trecut. Iată, toate lucrurile au devenit noi.

\VS{18}Dar toate lucrurile sunt de la Dumnezeu, care ne-a împăcat cu El însuși prin Isus Hristos și ne-a dat slujba împăcării;

\VS{19}și anume, că Dumnezeu a fost în Hristos împăcând lumea cu Sine, fără să le socotească greșelile lor, și ne-a încredințat nouă cuvântul împăcării.


\par }{\PP \VS{20}Noi suntem deci ambasadori în numele lui Hristos, ca și cum Dumnezeu ar ruga prin noi; vă rugăm în numele lui Hristos, împăcați-vă cu Dumnezeu.

\VS{21}Căci pe Cel ce n-a cunoscut nici un păcat, El L-a făcut să fie păcat în locul nostru, pentru ca noi să devenim în El neprihănirea lui Dumnezeu.



\par }\Chap{6}{\PP \VerseOne{1}Lucrând împreună, vă rugăm, de asemenea, să nu primiți în zadar harul lui Dumnezeu.

\VS{2}Căci el spune,

\par }{\Q “La un moment acceptabil te-am ascultat.

\par }{\Q În ziua mântuirii te-am ajutat.”

\par }{\PP Iată, acum este timpul potrivit. Iată, acum este ziua mântuirii.

\VS{3}Nu dăm prilej de poticnire în nimic, pentru ca slujba noastră să nu fie învinuită,

\VS{4}ci în orice lucru ne lăudăm ca slujitori ai lui Dumnezeu: în mare răbdare, în necazuri, în greutăți, în strâmtorări,

\VS{5}în bătăi, în închisori, în răscoale, în munci, în privegheri, în posturi,

\VS{6}în puritate, în cunoaștere, în perseverență, în bunătate, în Duhul Sfânt, în dragoste sinceră,

\VS{7}în cuvântul adevărului, în puterea lui Dumnezeu, prin armura dreptății la dreapta și la stânga,

\VS{8}prin glorie și dezonoare, prin raport rău și raport bun, ca înșelători și totuși adevărați,

\VS{9}ca necunoscuți și totuși bine cunoscuți, ca muribunzi și iată-ne vii, ca pedepsiți și nu uciși,

\VS{10}ca întristați și totuși mereu bucuroși, ca săraci și totuși îmbogățind pe mulți, ca neavând nimic și totuși posedând toate lucrurile.


\par }{\PP \VS{11}Gura noastră este deschisă pentru voi, Corinteni. Inima noastră este lărgită.

\VS{12}Voi nu sunteți limitați de noi, ci sunteți limitați de propriile voastre afecțiuni.

\VS{13}Acum, în schimb — vorbesc ca și copiilor mei — deschideți-vă și voi inimile.


\par }{\PP \VS{14}Nu vă legați cu necredincioșii printr-un jug nepotrivit, căci ce părtășie au dreptatea și nelegiuirea? Sau ce părtășie are lumina cu întunericul?

\VS{15}Ce înțelegere are Hristos cu Belial? Sau ce părtășie are un credincios cu un necredincios?

\VS{16}Ce înțelegere are un templu al lui Dumnezeu cu idolii? Căci voi sunteți un templu al Dumnezeului cel viu. Chiar așa cum a spus Dumnezeu: “Voi locui în ei și voi umbla în ei. Eu voi fi Dumnezeul lor și ei vor fi poporul meu”.

\VS{17}De aceea

\par }{\Q “'Ieși din mijlocul lor,

\par }{\QB și fiți despărțiți”, zice Domnul.

\par }{\Q 'Nu te atinge de nimic necurat.

\par }{\QB Te voi primi.


\par }{\Q \VS{18}Eu vă voi fi ca un Tată.

\par }{\QB Veți fi pentru mine fii și fiice”.

\par }{\PP spune Domnul Atotputernic.”



\par }\Chap{7}{\PP \VerseOne{1}Deci, iubiților, având aceste făgăduințe, să ne curățim de orice spurcăciune a cărnii și a duhului, desăvârșind sfințenia în frica lui Dumnezeu.


\par }{\PP \VS{2}Deschideți-ne inimile voastre. Nu am greșit față de nimeni. Nu am corupt pe nimeni. Nu am profitat de nimeni.

\VS{3}Nu spun aceasta ca să vă condamn, căci am mai spus că sunteți în inimile noastre pentru a muri împreună și a trăi împreună.

\VS{4}Mare este îndrăzneala mea de a vorbi față de voi. Mare este lăudăroșenia mea în numele vostru. Sunt plin de mângâiere. Debordez de bucurie în tot necazul nostru.


\par }{\PP \VS{5}Căci, când am ajuns în Macedonia, carnea noastră nu a avut nici o ușurare, ci am fost chinuiți din toate părțile. Luptele erau afară. Frica era înăuntru.

\VS{6}Cu toate acestea, Cel care îi mângâie pe cei umili, Dumnezeu, ne-a mângâiat prin venirea lui Titus,

\VS{7}și nu numai prin venirea lui, ci și prin mângâierea cu care a fost mângâiat în voi, în timp ce ne povestea despre dorul vostru, despre jalea voastră și despre zelul vostru pentru mine, așa că m-am bucurat și mai mult.


\par }{\PP \VS{8}Căci, deși v-am mâhnit cu scrisoarea mea, nu regret, deși am regretat-o. Căci văd că scrisoarea mea v-a întristat, deși numai pentru o vreme.

\VS{9}Acum mă bucur, nu pentru că v-ați întristat, ci pentru că v-ați întristat până la pocăință. Căci v-ați mâhnit în mod dumnezeiesc, ca să suferiți o pierdere din partea noastră în nimic.

\VS{10}Căci întristarea evlavioasă produce pocăința care duce la mântuire, care nu aduce niciun regret. Dar întristarea lumii produce moarte.

\VS{11}Căci iată, același lucru, că ați fost întristați în mod dumnezeiesc, ce grijă serioasă a lucrat în voi. Da, ce apărare, indignare, teamă, dorință, zel și răzbunare! În toate v-ați demonstrat că erați curați în această chestiune.

\VS{12}Așadar, deși v-am scris, nu am scris pentru cauza celui care a greșit, nici pentru cauza celui care a suferit greșeala, ci pentru ca grija voastră sinceră față de noi să se arate în voi în fața lui Dumnezeu.

\VS{13}De aceea am fost mângâiați. În mângâierea noastră, ne-am bucurat și mai mult pentru bucuria lui Titus, pentru că duhul lui a fost împrospătat de voi toți.

\VS{14}Căci, dacă în ceva m-am lăudat cu el în numele vostru, nu am fost dezamăgit. Ci, după cum v-am spus toate lucrurile în adevăr, tot așa și slava noastră, pe care am făcut-o înaintea lui Titus, s-a dovedit a fi adevăr.

\VS{15}Afecțiunea lui este și mai mare față de voi, în timp ce își amintește de toată ascultarea voastră, de modul în care l-ați primit cu frică și cu cutremur.

\VS{16}Mă bucur că în toate sunt încredințat cu privire la voi.



\par }\Chap{8}{\PP \VerseOne{1}Mai mult, fraților, vă facem cunoscut harul lui Dumnezeu, care a fost dat în adunările din Macedonia,

\VS{2}cum, într-o grea încercare de suferință, abundența bucuriei lor și sărăcia lor adâncă au sporit până la bogăția generozității lor.

\VS{3}Căci, după puterile lor, mărturisesc, da și dincolo de puterile lor, au dat de bunăvoie,

\VS{4}rugându-ne cu multă rugăminte să primim acest har și părtășia în slujirea sfinților.

\VS{5}Acest lucru nu a fost așa cum ne așteptam, ci mai întâi s-au dăruit pe ei înșiși Domnului și nouă prin voia lui Dumnezeu.

\VS{6}Așa că l-am îndemnat pe Titus, ca, după cum mai înainte făcuse un început, să desăvârșească și el în voi acest har.

\VS{7}Dar, așa cum voi abundați în toate — în credință, în exprimare, în cunoștință, în toată seriozitatea și în dragostea voastră față de noi —, vedeți să abundați și voi în acest har.


\par }{\PP \VS{8}Nu vorbesc ca o poruncă, ci ca să dovedesc, prin sinceritatea altora, sinceritatea dragostei voastre.

\VS{9}Căci cunoașteți harul Domnului nostru Isus Hristos, care, deși era bogat, s-a făcut sărac de dragul vostru, pentru ca voi, prin sărăcia lui, să vă îmbogățiți.

\VS{10}Eu vă dau un sfat în acest sens: se cuvine ca voi, care ați fost primii care ați început acum un an, nu numai să faceți, ci și să fiți binevoitori.

\VS{11}Dar acum săvârșiți și facerea, pentru ca, după cum a fost disponibilitatea de a fi dispuși, să fie și desăvârșirea din capacitatea voastră.

\VS{12}Căci, dacă disponibilitatea există, este acceptabil după ceea ce aveți, nu după ceea ce nu aveți.

\VS{13}Căci aceasta nu este pentru ca altora să le fie ușurat și vouă să vă chinuiți,

\VS{14}ci pentru egalitate. Abundența voastră în acest moment suplinește lipsa lor, pentru ca și abundența lor să devină o suplinire a lipsei voastre, pentru ca să fie egalitate.

\VS{15}După cum este scris: “Cel care a strâns mult nu a avut nimic de prisos, iar cel care a strâns puțin nu a avut lipsă”.


\par }{\PP \VS{16}Dar mulțumiri fie aduse lui Dumnezeu, care a pus în inima lui Titus aceeași grijă pentru voi.

\VS{17}Căci, într-adevăr, el a acceptat îndemnul nostru, dar, fiind el însuși foarte serios, a plecat de bunăvoie la voi.

\VS{18}Am trimis împreună cu el pe fratele a cărui laudă în Buna Vestire este cunoscută în toate adunările.

\VS{19}Și nu numai atât, ci și el a fost desemnat de adunări să călătorească împreună cu noi în acest har, pe care îl slujim pentru gloria Domnului însuși și pentru a ne arăta disponibilitatea noastră.

\VS{20}Noi evităm acest lucru, ca nimeni să ne învinuiască în legătură cu această abundență care este administrată de noi.

\VS{21}având în vedere lucruri onorabile, nu numai în ochii Domnului, ci și în ochii oamenilor.

\VS{22}Am trimis împreună cu ei pe fratele nostru, pe care l-am dovedit de multe ori harnic în multe lucruri, dar acum cu mult mai harnic, din cauza marii încrederi pe care o are în voi.

\VS{23}Cât despre Titus, el este partenerul și colaboratorul meu pentru voi. Cât despre frații noștri, ei sunt apostolii adunărilor, slava lui Hristos.

\VS{24}De aceea, arătați-le înaintea adunărilor dovada iubirii voastre față de ei și a lăudării noastre în numele vostru.



\par }\Chap{9}{\PP \VerseOne{1}Într-adevăr, este de prisos să vă scriu cu privire la slujba sfinților,

\VS{2}pentru că știu că sunteți pregătiți, și mă laud cu aceasta în numele vostru în fața celor din Macedonia, că și Ahaia a fost pregătită în anul trecut. Zelul vostru a stârnit foarte mulți dintre ei.

\VS{3}Dar am trimis pe frații aceștia pentru ca lăudarea noastră în numele vostru să nu fie în zadar în această privință, pentru ca, așa cum am spus, să fiți pregătiți,

\VS{4}ca nu cumva, dacă cineva din Macedonia vine acolo cu mine și vă găsește nepregătiți, noi (ca să nu mai vorbim de voi) să fim dezamăgiți în această lăudăroșenie încrezătoare.

\VS{5}De aceea, am considerat necesar să rog pe frați să meargă înainte la tine și să aranjeze din timp darul generos pe care l-ai promis mai înainte, pentru ca acesta să fie pregătit ca o chestiune de generozitate și nu de lăcomie.


\par }{\PP \VS{6}Amintește-ți că cine seamănă cu zgârcenie va secera și el cu zgârcenie. Cine seamănă din belșug va secera și el din belșug.

\VS{7}Fiecare să dea după cum a hotărât în inima sa, nu cu părere de rău sau cu silă, căci Dumnezeu iubește pe cel ce dă cu bucurie.

\VS{8}Și Dumnezeu poate să vă facă să vă prisosească tot harul, pentru ca, având totdeauna în toate cele de trebuință, să abundați pentru orice lucrare bună.

\VS{9}După cum este scris,

\par }{\Q “S-a împrăștiat în străinătate. El a dat săracilor.

\par }{\QB Neprihănirea Lui rămâne pentru totdeauna.”


\par }{\PP \VS{10}Iar Acela care dă sămânță semănătorului și pâine pentru hrană, să vă dea și El sămânța voastră pentru semănat și să înmulțească roadele dreptății voastre,

\VS{11}fiind voi îmbogățiți în toate pentru toată generozitatea, care aduce mulțumiri lui Dumnezeu prin noi.

\VS{12}Căci această slujbă de dăruire pe care o îndepliniți nu numai că suplinește lipsurile printre sfinți, ci abundă și prin multă mulțumire lui Dumnezeu,

\VS{13}având în vedere că, prin dovada dată de această slujbă, ei îl slăvesc pe Dumnezeu pentru ascultarea mărturisirii voastre față de Buna Vestire a lui Cristos și pentru generozitatea contribuției voastre față de ei și față de toți,

\VS{14}în timp ce ei înșiși, de asemenea, cu rugăminți în favoarea voastră, tânjesc după voi, datorită harului extraordinar al lui Dumnezeu în voi.

\VS{15}Acum, mulțumiri fie aduse lui Dumnezeu pentru darul său inefabil!



\par }\Chap{10}{\PP \VerseOne{1}Eu, Pavel, vă rog cu smerenia și blândețea lui Hristos, eu, care, în prezența voastră, sunt smerit între voi, dar, lipsind, sunt îndrăzneț față de voi.

\VS{2}Da, vă rog ca, atunci când sunt prezent, să nu-mi arăt curajul cu încrederea cu care intenționez să fiu îndrăzneț împotriva unora, care ne consideră că umblăm după trup.

\VS{3}Căci, deși umblăm în carne, nu ne luptăm după carne;

\VS{4}căci armele luptei noastre nu sunt din carne, ci puternice înaintea lui Dumnezeu pentru dărâmarea întăriturilor,

\VS{5}dărâmând închipuirile și orice înălțime care se ridică împotriva cunoașterii lui Dumnezeu și aducând orice gând în captivitate pentru supunerea lui Hristos,

\VS{6}și fiind gata să răzbune orice neascultare, când supunerea voastră va fi deplină.


\par }{\PP \VS{7}Privești lucrurile doar așa cum apar în fața ta? Dacă cineva se încrede în sine însuși că este al lui Hristos, să se gândească iarăși cu sine însuși că, după cum el este al lui Hristos, tot așa și noi suntem ai lui Hristos.

\VS{8}Căci chiar dacă mă laud oarecum din belșug cu privire la autoritatea noastră, pe care Domnul a dat-o pentru a vă zidi și nu pentru a vă doborî, nu mă voi rușina,

\VS{9}ca să nu pară că vreau să vă îngrozesc prin scrisorile mele.

\VS{10}Căci “scrisorile lui”, spun ei, “sunt grele și puternice, dar prezența lui trupească este slabă, iar vorbirea lui este disprețuită”.

\VS{11}O astfel de persoană să se gândească la acest lucru: ceea ce suntem prin scrisori în cuvânt când suntem absenți, așa suntem și în fapte când suntem prezenți.


\par }{\PP \VS{12}Căci nu îndrăznim să ne numărăm sau să ne comparăm cu unii dintre cei ce se recomandă. Dar ei înșiși, măsurându-se pe ei înșiși și comparându-se cu ei înșiși, sunt fără pricepere.

\VS{13}Noi însă nu ne vom lăuda dincolo de limitele cuvenite, ci în limitele cu care ne-a rânduit Dumnezeu, care ajung până la voi.

\VS{14}Căci nu ne întindem prea mult, ca și cum nu am ajunge până la voi. Căci am ajuns până la voi cu Vestea cea Bună a lui Hristos,

\VS{15}fără să ne lăudăm dincolo de limitele cuvenite, în lucrarea altora, ci având nădejdea că, pe măsură ce crește credința voastră, vom fi măriți din belșug de voi în sfera noastră de influență,

\VS{16}astfel încât să propovăduim Vestea cea Bună chiar și în părțile de dincolo de voi, nu să ne lăudăm cu ceea ce a făcut deja altcineva.

\VS{17}ci “cel care se laudă, să se laude în Domnul”.

\VS{18}Căci nu cel care se laudă pe sine este aprobat, ci cel pe care îl laudă Domnul.



\par }\Chap{11}{\PP \VerseOne{1}Aș vrea să mă îndurați cu puțină nebunie, dar, într-adevăr, mă îndurați.

\VS{2}Căci sunt gelos pe voi cu o gelozie dumnezeiască. Căci te-am făgăduit în căsătorie unui singur soț, ca să te prezint lui Hristos ca o fecioară curată.

\VS{3}Dar mă tem că, cumva, după cum șarpele a înșelat-o pe Eva cu viclenia lui, tot așa și mințile voastre ar putea fi corupte de la simplitatea care este în Hristos.

\VS{4}Căci, dacă cel care vine predică un alt Isus pe care noi nu l-am predicat, sau dacă primiți un alt spirit pe care nu l-ați primit, sau o altă “veste bună” pe care nu ați acceptat-o, ați suportat destul de bine acest lucru.

\VS{5}Căci socotesc că nu sunt deloc în urma celor mai buni apostoli.

\VS{6}Dar, chiar dacă sunt nepriceput în vorbire, nu sunt nepriceput în cunoaștere. Nu, în toate felurile vi s-a descoperit totul.


\par }{\PP \VS{7}Sau am păcătuit eu, umilindu-mă pe mine însumi, ca să vă înalț pe voi, pentru că v-am propovăduit gratuit vestea cea bună a lui Dumnezeu?

\VS{8}Am jefuit alte adunări, luând de la ele salariu pentru a vă sluji vouă.

\VS{9}Când eram prezent la voi și aveam nevoie, nu am fost o povară pentru nimeni, pentru că frații, când au venit din Macedonia, au suplinit măsura nevoilor mele. În toate m-am ferit să fiu o povară pentru voi și voi continua să fac acest lucru.

\VS{10}Cum adevărul lui Hristos este în mine, nimeni nu mă va opri din această laudă în regiunile din Ahaia.

\VS{11}De ce? Pentru că nu vă iubesc? Dumnezeu știe.


\par }{\PP \VS{12}Dar ceea ce fac eu, voi continua să fac, ca să tai prilejul celor ce doresc un prilej, ca să se laude, ca să fie recunoscuți ca și noi.

\VS{13}Căci astfel de oameni sunt falși apostoli, lucrători înșelători, care se dau drept apostoli ai lui Hristos.

\VS{14}Și nu este de mirare, căci chiar și Satana se deghizează în înger de lumină.

\VS{15}Prin urmare, nu este mare lucru dacă și slujitorii lui se dau drept slujitori ai dreptății, al căror sfârșit va fi potrivit cu faptele lor.


\par }{\PP \VS{16}Și iarăși zic: să nu mă creadă nimeni nebun. Ci, dacă este așa, primiți-mă totuși ca pe un nebun, ca să mă laud și eu puțin.

\VS{17}Ceea ce vorbesc, nu vorbesc după Domnul, ci ca o nebunie, în această încredere de laudă.

\VS{18}Văzând că mulți se laudă după trup, mă voi lăuda și eu.

\VS{19}Căci voi suportați cu plăcere pe cei nebuni, fiind înțelepți.

\VS{20}Căci voi suportați pe un om dacă vă duce în robie, dacă vă devorează, dacă vă ia în captivitate, dacă se înalță sau dacă vă lovește peste față.

\VS{21}Spre rușinea mea, vorbesc ca și cum am fi fost slabi. Totuși, în orice fel în care cineva este îndrăzneț (vorbesc în neștire), sunt și eu îndrăzneț.

\VS{22}Sunt ei evrei? La fel și eu, sunt ei israeliți? La fel sunt și eu. Sunt ei urmașii lui Avraam? La fel sunt și eu.

\VS{23}Sunt ei slujitori ai lui Hristos? (Vorbesc ca unul de lângă el însuși.) Eu sunt mai mult: în munci mai mult, în închisori mai mult, în bătăi peste măsură și în morți des.

\VS{24}De cinci ori am primit de la iudei patruzeci de lovituri minus una.

\VS{25}De trei ori am fost bătut cu toiege. O dată am fost ucis cu pietre. De trei ori am suferit un naufragiu. Am stat o noapte și o zi în adânc.

\VS{26}Am fost deseori în călătorii, primejdii de râuri, primejdii de tâlhari, primejdii de la compatrioții mei, primejdii de la neamuri, primejdii în cetate, primejdii în pustiu, primejdii pe mare, primejdii printre frații mincinoși;

\VS{27}în munci și osteneli, în privegheri dese, în foame și sete, în posturi dese, în frig și goliciune.


\par }{\PP \VS{28}În afară de cele de afară, mai este ceva care mă apasă în fiecare zi: îngrijorarea pentru toate adunările.

\VS{29}Cine este slab, și eu nu sunt slab? Cine este făcut să se poticnească, și eu nu ard de indignare?


\par }{\PP \VS{30}Dacă trebuie să mă laud, mă voi lăuda cu ceea ce ține de slăbiciunea mea.

\VS{31}Dumnezeul și Tatăl Domnului Isus Hristos, cel binecuvântat în veci de veci, știe că nu mint.

\VS{32}În Damasc, guvernatorul subordonat regelui Aretas a păzit cetatea Damascului, dorind să mă aresteze.

\VS{33}Am fost coborât într-un coș pe o fereastră lângă zid și am scăpat din mâinile lui.



\par }\Chap{12}{\PP \VerseOne{1}Fără îndoială că nu este de folos să mă laud, dar voi ajunge la viziuni și revelații ale Domnului.

\VS{2}Cunosc un om în Hristos care a fost răpit în al treilea cer acum paisprezece ani — dacă în trup, nu știu, sau dacă în afara trupului, nu știu; Dumnezeu știe.

\VS{3}Cunosc un astfel de om (fie că era în trup, fie că era în afara trupului, nu știu; Dumnezeu știe),

\VS{4}cum a fost răpit în Paradis și a auzit cuvinte de nedescris, pe care nu-i este permis unui om să le rostească.

\VS{5}În numele unui astfel de om mă voi lăuda, dar în numele meu nu mă voi lăuda, decât în slăbiciunile mele.

\VS{6}Căci, dacă aș vrea să mă laud, nu voi fi nebun, căci voi spune adevărul. Dar mă abțin, pentru ca nimeni să nu creadă despre mine mai mult decât ceea ce vede în mine sau aude de la mine.

\VS{7}Din cauza mărimii prea mari a revelațiilor, ca să nu mă înalț prea mult, mi s-a dat un spin în carne: un mesager al lui Satana care să mă chinuie, ca să nu mă înalț prea mult.

\VS{8}În legătură cu acest lucru, l-am rugat de trei ori pe Domnul ca să se îndepărteze de la mine.

\VS{9}El mi-a spus:
{\WJ{“Harul meu îți este de ajuns, căci puterea mea se desăvârșește în slăbiciune”. De aceea, cu cea mai mare plăcere mă voi lăuda mai degrabă în slăbiciunile mele, pentru ca puterea lui Hristos să se odihnească peste mine.
}}


\par }{\PP \VS{10}De aceea îmi plac slăbiciunile, vătămările, necazurile, nevoile, persecuțiile și strâmtorările, pentru Hristos. Căci atunci când sunt slab, atunci sunt puternic.

\VS{11}Am devenit nebun în a mă lăuda. M-ați constrâns, pentru că ar fi trebuit să fiu lăudat de voi, căci nu sunt cu nimic mai prejos decât cei mai buni apostoli, deși nu sunt nimic.

\VS{12}Cu adevărat, semnele unui apostol s-au făcut printre voi cu toată perseverența, prin semne, minuni și lucrări puternice.

\VS{13}Căci în ce anume ați fost făcuți inferiori celorlalte adunări, dacă nu în faptul că eu însumi nu am fost o povară pentru voi? Iertați-mi această greșeală!


\par }{\PP \VS{14}Iată, este a treia oară când sunt gata să vin la voi și nu vreau să fiu o povară pentru voi, căci nu pe voi vă caut averile voastre, ci pe voi. Căci nu copiii trebuie să economisească pentru părinți, ci părinții pentru copii.

\VS{15}Cu mare plăcere voi cheltui și voi fi cheltuit pentru sufletele voastre. Dacă vă iubesc mai mult și mai mult, sunt eu oare mai puțin iubit?

\VS{16}Chiar și așa, eu însumi nu v-am împovărat. Dar ați putea spune că, fiind viclean, v-am prins cu înșelăciune.

\VS{17}Oare am profitat de voi prin vreunul dintre cei pe care vi i-am trimis?

\VS{18}L-am îndemnat pe Titus și am trimis împreună cu el pe fratele său. A profitat oare Titus de voi? Nu cumva am umblat în același spirit? Nu am mers pe aceiași pași?


\par }{\PP \VS{19}Iarăși, credeți că ne scuzăm în fața voastră? În fața lui Dumnezeu, noi vorbim în Hristos. Dar toate lucrurile, iubiților, sunt pentru zidirea voastră.

\VS{20}Căci mă tem că poate, atunci când voi veni, nu vă voi găsi așa cum vreau eu și că aș putea fi găsit de voi așa cum nu doriți voi, că poate vor fi certuri, gelozii, izbucniri de mânie, facțiuni, calomnii, șoapte, gânduri de mândrie sau revolte,

\VS{21}că, din nou, când voi veni, Dumnezeul meu mă va umili înaintea voastră și voi plânge pentru mulți dintre cei care au păcătuit până acum și nu s-au pocăit de necurăția, imoralitatea sexuală și poftele pe care le-au săvârșit.



\par }\Chap{13}{\PP \VerseOne{1}Este a treia oară când vin la voi. “În gura a doi sau trei martori se va adeveri orice cuvânt.”

\VS{2}Am avertizat mai înainte și avertizez din nou, așa cum am fost prezent a doua oară, tot așa și acum, fiind absent, scriu celor care au păcătuit mai înainte acum și tuturor celorlalți că, dacă voi veni din nou, nu voi cruța,

\VS{3}având în vedere că voi căutați o dovadă a lui Cristos care vorbește în mine, care nu este slab, ci este puternic în voi.

\VS{4}Căci el a fost răstignit prin slăbiciune, dar trăiește prin puterea lui Dumnezeu. Căci și noi suntem slabi în el, dar vom trăi împreună cu el prin puterea lui Dumnezeu față de voi.


\par }{\PP \VS{5}Cercetați-vă pe voi înșivă dacă sunteți în credință. Testați-vă pe voi înșivă. Sau nu știți despre voi înșivă că Isus Hristos este în voi?” — dacă nu cumva sunteți într-adevăr descalificați.

\VS{6}Dar sper că veți ști că nu suntem descalificați.


\par }{\PP \VS{7}Și mă rog lui Dumnezeu să nu faceți nimic rău, nu ca noi să fim învredniciți, ci ca voi să faceți ceea ce este onorabil, chiar dacă pare că am dat greș.

\VS{8}Căci nu putem face nimic împotriva adevărului, ci pentru adevăr.

\VS{9}Căci ne bucurăm când noi suntem slabi și voi sunteți puternici. Ne rugăm și pentru aceasta: pentru ca voi să deveniți desăvârșiți.

\VS{10}De aceea, scriu aceste lucruri în absență, ca să nu mă port cu asprime când sunt prezent, conform autorității pe care mi-a dat-o Domnul pentru zidire și nu pentru dărâmare.


\par }{\PP \VS{11}În sfârșit, fraților, bucurați-vă! Fiți desăvârșiți. Fiți mângâiați. Fiți de aceeași părere. Trăiți în pace, și Dumnezeul iubirii și al păcii va fi cu voi.

\VS{12}Salutați-vă unii pe alții cu o sărutare sfântă.


\par }{\PP \VS{13}Toți sfinții vă salută.


\par }{\PP \VS{14}Harul Domnului Isus Hristos, dragostea lui Dumnezeu și părtășia Duhului Sfânt să fie cu voi toți. Amin.


\par }